%!TeX root=../main.tex

\chapter{Sobre o kworkflow}

\label{cap:kworkflow}

O kw só se entende a partir do trabalho que ele automatiza, por isso este
capítulo começa pelo desenvolvimento do kernel Linux. Depois apresenta a
ferramenta e termina no dado que o resto do trabalho usa: o registro que o kw
faz do próprio uso e a base local onde ele fica.

\section{O modelo de desenvolvimento do kernel Linux}
\label{sec:modelo-kernel}

O kernel Linux é desenvolvido por milhares de pessoas ao mesmo tempo. Na
versão 7.3, lançada em outubro de 2026, 3.021 desenvolvedores contribuíram
com mais de 18 mil \textit{commits}, que acrescentaram cerca de 650 mil linhas
líquidas ao código; 807 dessas pessoas nunca haviam contribuído antes
\citep{corbet2026stats73}. Nenhuma pessoa consegue revisar esse volume
sozinha, e por isso o projeto é dividido em subsistemas. Cada subsistema
cuida de uma parte do código, como a gerência de memória, a pilha de rede ou
um conjunto de \textit{drivers}, e tem os próprios mantenedores, a própria
árvore \textit{git} e a própria lista de e-mail \citep{kerneldocprocess}. 

Uma alteração sobe por uma cadeia de confiança até a \textit{mainline}, a
árvore mantida por Linus Torvalds e da qual saem as versões oficiais
(Figura~\ref{fig:caminho-patch}). O desenvolvedor envia o trabalho ao
mantenedor do \textit{driver} ou do arquivo que modificou; este o repassa ao
mantenedor do subsistema; e o mantenedor do subsistema, no momento certo do
ciclo, pede a Torvalds que integre a sua árvore à \textit{mainline}, por meio
de um \textit{pull request}\footnote{Aqui, \textit{pull request} é o pedido
para que uma árvore \textit{git} seja integrada a outra, e não o recurso de
mesmo nome do GitHub.}. Torvalds não tem como revisar cada alteração e confia
que os níveis de baixo já fizeram esse trabalho.

\begin{figure}[htbp]
  \centering
  \resizebox{\textwidth}{!}{%!TeX root=../main.tex
% Figura 2.1: caminho de um patch até a mainline.
% Adaptada de Kroah-Hartman (2018), images/kernel_patch_process.svg em
% github.com/gregkh/kernel-development, CC BY-SA 4.0. Mantém a disposição e
% as arestas do original; acrescenta a faixa da interface
% contribuidor-mantenedor e traduz os rótulos.
\begin{tikzpicture}[
    x=1cm, y=1cm,
    every node/.style={font=\footnotesize, align=center},
    no/.style={draw, rounded corners=4pt, minimum width=2.1cm,
               minimum height=0.8cm, inner sep=2pt},
    topo/.style={no, fill=black!12},
    subs/.style={no, fill=black!6},
    mant/.style={no, fill=white},
    dev/.style={no, fill=white, minimum width=1.6cm, dashed, font=\scriptsize},
    seta/.style={-{Latex[length=1.6mm]}, semithick},
  ]
  % Topo
  \node[topo] (next)   at (5.4, 6.1) {árvore \texttt{-next}};
  \node[topo] (linus)  at (2.6, 5.0) {Linus Torvalds\\(mainline)};
  \node[topo] (morton) at (8.2, 5.0) {Andrew Morton};

  % Mantenedores de subsistema
  \node[subs] (s1) at (1.9, 3.5) {mantenedor de\\subsistema};
  \node[subs] (s2) at (4.6, 3.5) {mantenedor de\\subsistema};
  \node[subs] (s3) at (7.3, 3.5) {mantenedor de\\subsistema};

  % Mantenedores de driver ou arquivo
  \node[mant] (m1) at (0.9, 2.0) {mantenedor de\\driver/arquivo};
  \node[mant] (m2) at (3.6, 2.0) {mantenedor de\\driver/arquivo};
  \node[mant] (m3) at (6.3, 2.0) {mantenedor de\\driver/arquivo};
  \node[mant] (m4) at (9.0, 2.0) {mantenedor de\\driver/arquivo};

  % Desenvolvedores
  \foreach \i/\x in {1/0.0, 2/1.9, 3/3.8, 4/5.7, 5/7.6, 6/9.5} {
    \node[dev] (d\i) at (\x, 0) {desenvolvedor};
  }

  % Faixa da interface contribuidor-mantenedor (acréscimo desta adaptação)
  \draw[thick, dashed, rounded corners=6pt]
    (-1.0, 0.62) rectangle (10.5, 1.18);

  % Desenvolvedor -> primeiro mantenedor
  \draw[seta] (d1) -- (m1);
  \draw[seta] (d2) -- (m1);
  \draw[seta] (d2) -- (m2);
  \draw[seta] (d4) -- (m3);
  \draw[seta] ([xshift=3mm]d4.north) -- (7.65, 2.0) -- ([xshift=2mm]s3.south);
  \draw[seta] (d3) -- (m3);
  \draw[seta] (d5) -- (m4);
  \draw[seta] (d6.north) to[out=90, in=-90] ([xshift=0.35cm]m4.east) to[out=90, in=-55] (morton);

  % Mantenedor de driver/arquivo -> subsistema
  \draw[seta] (m1) -- (s1);
  \draw[seta] (m2) -- (s1);
  \draw[seta] (m3) -- (s2);
  \draw[seta] (m4) -- (morton);

  % Subsistema -> Linus e -next
  \draw[seta] (s1) -- (linus);
  \draw[seta] (s2) -- (linus);
  \draw[seta] (s3) -- (linus);
  \draw[seta] (s1) -- (next);
  \draw[seta] (s2) -- (next);
  \draw[seta] (s3) -- (next);

  % Andrew Morton -> Linus e -next
  \draw[seta] (morton) -- (linus);
  \draw[seta] (morton) -- (next);
  % Rótulo da faixa por cima das setas
  \node[font=\footnotesize\itshape, fill=white, inner sep=1.5pt]
    at (4.75, 0.9) {interface contribuidor--mantenedor};
\end{tikzpicture}
}
  \caption{Caminho de uma alteração até a \textit{mainline}. Adaptado de
    \citet{kroah2018kerneldev}, sob licença CC BY-SA~4.0; a faixa da
    interface contribuidor--mantenedor foi acrescentada nesta adaptação.}
  \label{fig:caminho-patch}
\end{figure}

A contribuição em si chega por e-mail. A unidade de trabalho é o
\textit{patch}: um arquivo de texto com a diferença entre duas versões do
código e uma mensagem que explica o que mudou e por quê. Parece um
\textit{commit} do \textit{git}, mas formatado para ser enviado a uma lista de
discussão \citep{tadokoro2025sustain}. É na lista que a revisão acontece. Os
mantenedores e outros desenvolvedores respondem ao e-mail com comentários, o
autor corrige e reenvia uma nova versão, e o ciclo se repete até que o
\textit{patch} seja aceito ou abandonado \citep{kerneldocprocess}. Toda essa
troca fica pública e arquivada no
\textit{lore.kernel.org}\footnote{\url{https://lore.kernel.org}}, onde
qualquer pessoa pode ler a discussão de um \textit{patch} do envio ao aceite
\citep{lorekernel, tadokoro2023lore}.

O ritmo das versões é quase fixo: cada ciclo abre com uma janela de cerca de
duas semanas em que os mantenedores enviam a Torvalds o que acumularam, e
segue só com correções até a versão final \citep{kerneldocprocess}. Os
subsistemas aceitam \textit{patches} o tempo todo, mas cada um só chega à
\textit{mainline} no ritmo dos ciclos.

O modelo funciona há mais de trinta anos, mas tem custos. Entre as versões
3.19 (2015) e 6.13 (2025), o código e o número de mantenedores cresceram de
forma constante, enquanto o volume de \textit{commits} por ciclo oscilou sem
tendência de alta \citep{tadokoro2025sustain}. Número de \textit{commits} não
mede esforço, advertem os autores, mas o dado sugere que mais mantenedores não
se traduziram em mais vazão. Parte das contribuições nem chega a ser revisada:
no estudo de \citet{rigby2014peer} sobre as listas do kernel entre 2005 e
2008, 44\% dos \textit{patches} enviados não receberam discussão de revisão. O
problema não diminuiu com o tempo. Ao comentar a chegada recorde de novos
desenvolvedores na versão 7.3, \citet{corbet2026stats73} escreve que os
mantenedores já estão perdendo a capacidade de acompanhar o fluxo. A conclusão
de \citet{corbet2017lfreport} continua valendo: sem as ferramentas certas, um
projeto do tamanho do kernel não conseguiria funcionar sem desabar sob o
próprio peso.

\begin{figure}[htbp]
  \centering
  \resizebox{\textwidth}{!}{%!TeX root=../main.tex
% Figura 2.2: fluxo de desenvolvimento de um patchset.
% Adaptada de Tadokoro, Siqueira e Meirelles (2025), Fig. 1 ("Concept map of
% the patchset development workflow"). Mantém as etapas e os laços de falha;
% traduz os rótulos e destaca as duas etapas que o kw registra (compilação e
% instalação), retomadas na seção 2.4.
\begin{tikzpicture}[
    x=1cm, y=1cm,
    every node/.style={font=\footnotesize, align=center},
    etapa/.style={draw, rounded corners=3pt, minimum width=2.5cm,
                  minimum height=0.95cm, inner sep=2pt, fill=white},
    reg/.style={etapa, very thick, fill=black!12},
    borda/.style={etapa, rounded corners=9pt, fill=black!4},
    seta/.style={-{Latex[length=1.6mm]}, semithick},
    falha/.style={-{Latex[length=1.6mm]}, semithick, dashed},
    rot/.style={font=\scriptsize\itshape, fill=white, inner sep=1pt},
  ]
  % Linha de cima: da concepção à instalação
  \node[borda] (conc) at (0.0, 2.6) {concepção do\\patchset};
  \node[etapa] (impl) at (3.4, 2.6) {implementação};
  \node[reg]   (comp) at (6.8, 2.6) {compilação\\(\textit{build})};
  \node[reg]   (inst) at (10.2, 2.6) {instalação\\(\textit{deploy})};

  % Linha de baixo: dos testes ao merge
  \node[etapa] (test) at (10.2, 0.0) {testes\\(automáticos e manuais)};
  \node[etapa] (envi) at (6.8, 0.0) {envio aos\\mantenedores e listas};
  \node[etapa] (revi) at (3.4, 0.0) {revisão};
  \node[borda] (merg) at (0.0, 0.0) {integração\\à árvore};

  % Caminho de sucesso
  \draw[seta] (conc) -- (impl);
  \draw[seta] (impl) -- (comp);
  \draw[seta] (comp) -- (inst);
  \draw[seta] (inst) -- (test);
  \draw[seta] (test) -- (envi);
  \draw[seta] (envi) -- (revi);
  \draw[seta] (revi) -- (merg);

  % Laços de falha, todos de volta à implementação
  \draw[falha] (comp.north) -- ++(0, 0.5) -| ([xshift=-4mm]impl.north);
  \node[rot, anchor=south] at (5.1, 3.55) {falha na compilação};
  \draw[falha] (inst.north) -- ++(0, 1.1) -| ([xshift=4mm]impl.north);
  \node[rot, anchor=south] at (6.8, 4.15) {falha na instalação};
  \draw[falha] (test.south) |- (1.8, -0.95) |- ([yshift=-2.5mm]impl.west);
  \node[rot, anchor=north] at (6.8, -1.0) {falha nos testes};
  \draw[falha] (revi.north) -- (impl.south)
    node[rot, midway, right=2pt] {alterações pedidas};
\end{tikzpicture}
}
  \caption{Fluxo de desenvolvimento de um conjunto de \textit{patches}, do
    ponto de vista de quem contribui. As etapas em destaque são as que o
    kworkflow registra, retomadas na Seção~\ref{sec:base-local}. Adaptado de
    \citet{tadokoro2025kw}.}
  \label{fig:fluxo-patchset}
\end{figure}

Uma parte desse custo cai sobre o desenvolvedor antes mesmo de o
\textit{patch} chegar à lista (Figura~\ref{fig:fluxo-patchset}). Ele precisa
compilar o kernel, instalá-lo numa máquina de teste, verificar se a máquina
inicializa e se os testes passam, e só então preparar e enviar o
\textit{patch}. Cada falha o leva de volta à implementação. Nenhuma dessas
etapas é trivial no kernel: uma compilação completa pode levar de minutos a
horas, e a máquina de teste pode ser virtual, remota ou o próprio computador
do desenvolvedor \citep{tadokoro2025kw}. Para cada etapa existe uma
ferramenta: o \texttt{make}, os \textit{scripts} \texttt{checkpatch.pl} e
\texttt{get\_maintainer.pl} que acompanham o código-fonte, o \texttt{git
send-email}, o QEMU. Elas resolvem cada uma a sua parte, e juntá-las num fluxo
de trabalho fica a cargo de cada desenvolvedor, em geral com \textit{scripts}
pessoais que ninguém mais usa \citep{tadokoro2025kw}. O kworkflow nasceu para
reunir esse trabalho num lugar só.

\section{O que é e a que se propõe}
\label{sec:kw-o-que-e}

O kworkflow (kw) é uma ferramenta de software livre, distribuída sob a GPL
versão~2 e escrita em Bash, que reúne numa só interface de linha de comando as
tarefas descritas na seção anterior. O projeto resume o propósito numa frase:
reduzir o custo de preparação do ambiente para quem trabalha com o kernel
Linux e oferecer ferramentas para as tarefas do dia a dia
\citep{kworkflowrepo}. \citet{tadokoro2025kw} o classificam como um
\textit{Developer Automation Workflow System} (DAWS), termo que
\citet{tadokoro2023lore} propôs com a comunidade do projeto por analogia a
\textit{Version Control System}: assim como um sistema de controle de versão
cuida do histórico do código, um DAWS cuida do fluxo de trabalho em volta
dele.

Todos os comandos seguem a forma \texttt{kw <comando> <opções>}, e o projeto
tem hoje duas dezenas deles \citep{tadokoro2025kw}. Eles se dividem em dois
grupos \citep{tadokoro2023lore}. Os comandos práticos agem diretamente sobre o
kernel, como \texttt{kw build}, que compila, e \texttt{kw deploy}, que instala
o kernel compilado numa máquina de teste. Os comandos indiretos não mexem no
kernel, mas melhoram o trabalho em volta dele: gerenciar máquinas remotas,
enviar \textit{patches}, consultar o arquivo das listas, cronometrar sessões
de trabalho. A Tabela~\ref{tab:comandos-kw} reúne alguns comandos de cada
grupo. Para ter uma ideia do que o kw economiza, compilar o
kernel para um processador ARM a partir de um computador x86, com todos os
núcleos e uma \textit{flag} de compilação, exige uma linha de \texttt{make}
com quatro parâmetros; depois de configurado uma vez, o kw reduz isso a
\texttt{kw build} \citep{tadokoro2025kw}.

\begin{table}[htbp]
  \centering
  \caption{Alguns comandos do kworkflow. Seleção a partir da lista completa de
    \citet{tadokoro2025kw}.}
  \label{tab:comandos-kw}
  \begin{tabular}{@{}lll@{}}
    \toprule
    Comando & Tipo & O que faz \\
    \midrule
    \texttt{build}                 & prático  & compila o kernel e os módulos \\
    \texttt{deploy}                & prático  & instala o kernel numa máquina de teste \\
    \texttt{kernel-config-manager} & prático  & gerencia arquivos \texttt{.config} \\
    \texttt{send-patch}            & indireto & envia \textit{patches} por e-mail \\
    \texttt{patch-hub}             & indireto & navega pelos \textit{patches} do \textit{lore.kernel.org} \\
    \texttt{pomodoro}              & indireto & cronometra sessões de trabalho \\
    \texttt{report}                & indireto & mostra estatísticas de uso \\
    \bottomrule
  \end{tabular}
\end{table}


Por dentro, o kw é organizado como um \textit{hub}
(Figura~\ref{fig:arquitetura-kw}). Um único ponto de entrada, o arquivo
\texttt{kw} na raiz do repositório, recebe o comando e o despacha para o
componente correspondente, como \texttt{src/build.sh} para \texttt{kw build}.
Os componentes compartilham bibliotecas em \texttt{src/lib}, que tratam de
tarefas comuns como manipular texto, acessar máquinas remotas ou ler e gravar
no banco de dados. O código que depende de uma distribuição, de uma
arquitetura ou de um subsistema específico, e que por isso muda com
frequência, fica isolado em \textit{plugins} \citep{tadokoro2025kw,
neto2022kw}. A mesma ideia de \textit{hub} vale para a relação com outras
ferramentas: o kw prefere incorporar o que a comunidade já usa, como o
\texttt{git}, o \texttt{checkpatch.pl} e o QEMU, e só escreve uma solução
própria quando nada resolve o problema \citep{tadokoro2025kw}. Isso evita mais
uma camada de \textit{scripts} duplicados, que é justamente o problema
descrito na Seção~\ref{sec:modelo-kernel}.

\begin{figure}[htbp]
  \centering
  \resizebox{0.85\textwidth}{!}{%!TeX root=../main.tex
% Figura 2.3: arquitetura conceitual do kworkflow.
% Adaptada de Tadokoro, Siqueira e Meirelles (2025), Fig. 6 ("kw conceptual
% architecture"). Mantém as camadas do original e destaca as duas
% bibliotecas que gravam os dados de uso (kwlib, onde fica statistics_manager,
% e kw_db), retomadas na
% seção 2.4. Como no original, as ligações são ilustrativas.
\begin{tikzpicture}[
    x=1cm, y=1cm,
    every node/.style={font=\footnotesize, align=center},
    peca/.style={draw, rounded corners=3pt, minimum width=1.9cm,
                 minimum height=0.75cm, inner sep=2pt, fill=white,
                 font=\footnotesize\ttfamily},
    dado/.style={peca, very thick, fill=black!12},
    camada/.style={draw, dashed, rounded corners=6pt, inner sep=6pt},
    rotulo/.style={font=\scriptsize\scshape, anchor=west},
    seta/.style={-{Latex[length=1.6mm]}, semithick},
  ]
  % Hub
  \node[peca, minimum width=2.6cm, fill=black!6] (hub) at (4.5, 3.6)
    {kw\\[-1pt]{\normalfont\scriptsize(ponto de entrada)}};

  % Componentes
  \node[peca] (c1) at (0.8, 2.0) {build};
  \node[peca] (c2) at (3.0, 2.0) {deploy};
  \node[peca] (c3) at (5.2, 2.0) {send-patch};
  \node[peca] (c4) at (7.4, 2.0) {report};
  \node (c5) at (9.0, 2.0) {\dots};

  % Bibliotecas
  \node[dado] (l1) at (0.8, 0.4) {kw\_db};
  \node[dado] (l2) at (3.0, 0.4) {kwlib};
  \node[peca] (l3) at (5.2, 0.4) {kw\_string};
  \node[peca] (l4) at (7.4, 0.4) {remote};
  \node (l5) at (9.0, 0.4) {\dots};

  % Plugins
  \node[peca, minimum width=2.7cm] (p1) at (2.6, -1.2) {kernel\_install};
  \node[peca, minimum width=2.3cm] (p2) at (5.4, -1.2) {subsystems};
  \node (p3) at (7.2, -1.2) {\dots};

  % Molduras das camadas
  \node[camada, fit=(c1)(c5)] (fc) {};
  \node[camada, fit=(l1)(l5)] (fl) {};
  \node[camada, fit=(p1)(p3)] (fp) {};
  \node[rotulo] at ([xshift=2pt, yshift=7pt]fc.north west) {componentes};
  \node[rotulo] at ([xshift=2pt, yshift=7pt]fl.north west) {bibliotecas};
  \node[rotulo] at ([xshift=2pt, yshift=7pt]fp.north west) {plugins};

  % O hub despacha para os componentes
  \foreach \c in {c1, c2, c3, c4} {
    \draw[seta] (hub.south) -- (\c.north);
  }
  % Os componentes usam bibliotecas e plugins (ligações ilustrativas)
  \draw[seta] (fc.south) -- (fl.north);
  \draw[seta] (fl.south) -- (fp.north);
\end{tikzpicture}
}
  \caption{Arquitetura conceitual do kworkflow. As bibliotecas em destaque
    são as que gravam os dados de uso, descritas na Seção~\ref{sec:base-local}.
    As ligações entre camadas são ilustrativas. Adaptado de
    \citet{tadokoro2025kw}.}
  \label{fig:arquitetura-kw}
\end{figure}

O projeto começou em maio de 2017, como um conjunto de \textit{scripts} para
compilar e instalar o kernel, e com o tempo se tornou uma ferramenta que cobre
o fluxo inteiro \citep{kworkflowrepo, neto2022kw}. O desenvolvimento acontece
no GitHub: as contribuições entram por \textit{pull request} na
\textit{branch} \texttt{unstable}, e a \textit{branch} \texttt{master} recebe
as versões estáveis \citep{neto2022kw}. Rodrigo Siqueira e Matheus Tavares
são os autores principais \citep{tadokoro2025kw}, e boa parte da evolução
recente veio de trabalhos de conclusão de curso do IME-USP.
\citet{neto2022kw} reorganizou a documentação, criou o primeiro modelo de
banco de dados e o comando de envio de \textit{patches};
\citet{tadokoro2023lore} concluiu a migração para o banco e criou o
\texttt{patch-hub}, que integra o kw ao \textit{lore.kernel.org}; e
\citet{souza2025kw} criou o acompanhamento dos \textit{patches} já enviados.
Cada um acrescentou ao kw uma funcionalidade nova. Este trabalho segue a mesma
linha. Propõe uma forma de levar para fora da máquina do desenvolvedor os dados
que o kw já guarda.

Esses dados não são acidentais. Os autores do kw apresentam a ferramenta
também como plataforma de pesquisa: segundo \citet{tadokoro2025kw}, a
infraestrutura de coleta de dados que o kw já tem o torna um bom ponto de
partida para estudar empiricamente o modelo de desenvolvimento do kernel. As
duas seções seguintes descrevem o que, onde e como o kw registra o uso hoje.

\section{O registro de uso e o kw report}
\label{sec:funcionalidades}

O kw registra o próprio uso sem que o usuário precise fazer nada. Ao fim de
cada compilação e de cada instalação, o componente correspondente chama uma
única função, \texttt{statistics\_manager}, com o nome da operação, o momento
em que ela começou, quanto tempo durou e se deu certo
(Programa~\ref{prog:registro}). Hoje só \texttt{build.sh} e
\texttt{deploy.sh} fazem essa chamada, com cinco rótulos: \texttt{build},
\texttt{deploy}, \texttt{modules\_deploy}, \texttt{list} e
\texttt{uninstall}. O banco já prevê rótulos para todos os outros comandos,
mas nenhum deles grava registro. Cada execução termina num de quatro
estados: \texttt{success}, \texttt{failure}, \texttt{interrupted} ou
\texttt{unknown}. A coleta vem ligada por padrão e pode ser desligada pela
opção \texttt{disable\_statistics\_data\_track} do arquivo de configuração
\citep{kworkflowrepo, tadokoro2025kw}. Nada disso sai da máquina.

\begin{program}[htbp]
  \centering
\begin{lstlisting}[language=bash, style=wider]
# src/build.sh, ao fim de kw build
statistics_manager 'build' "$start" "$runtime" 'failure' "$flag"  # falhou
statistics_manager 'build' "$start" "$runtime" '' "$flag"         # sucesso

# src/lib/kwlib.sh
function statistics_manager()
{
  local label_name="$1"
  local start_datetime_in_secs="$2"
  local elapsed_time_in_secs="$3"
  local status=${4:-'success'}
  ...
\end{lstlisting}
  \caption{Registro de uma compilação: as chamadas em \texttt{src/build.sh}
    e o início da função \texttt{statistics\_manager}, em
    \texttt{src/lib/kwlib.sh} \citep{kworkflowrepo}.}
  \label{prog:registro}
\end{program}

O único lugar onde esses registros são lidos é o \texttt{kw report}. Ele
agrega as execuções por dia, semana, mês ou ano e mostra, para cada
operação, quantas vezes ela rodou e o maior, o menor e o tempo médio de
execução (Programa~\ref{prog:report}). A coluna \enquote{Total} é o número de
execuções. O relatório não separa sucessos de
falhas: o estado fica guardado, mas uma compilação que falhou entra na mesma
contagem de uma que deu certo. O mesmo comando resume ainda as sessões do
cronômetro de foco do kw, que ficam fora do escopo deste trabalho. O
\texttt{kw report} roda na máquina de quem gerou os dados e mostra o
resultado só para essa pessoa.

\begin{program}[htbp]
  \centering
\begin{lstlisting}[style=wider, language={}, style=nocolor, numbers=none]
$ kw report --statistics --month
# Statistics Report: month 2026/10
               Total  Max      Min      Average
Build             21 00:23:25 00:03:07 00:12:18
Deploy            15 00:01:10 00:00:28 00:00:45
Modules_deploy     4 00:00:27 00:00:18 00:00:21
List               5 00:00:03 00:00:02 00:00:02
Uninstall          2 00:00:05 00:00:04 00:00:04
\end{lstlisting}
  \caption{Exemplo de saída do \texttt{kw report}. O formato é o do kw; os
    valores são ilustrativos.}
  \label{prog:report}
\end{program}

\section{A base local de estatísticas}
\label{sec:base-local}

Todos esses registros ficam num único arquivo SQLite, o \texttt{kw.db}, no
diretório de dados do usuário. O kw o acessa pela linha de comando do
\texttt{sqlite3}, sem precisar de outra linguagem além do Bash, e esse foi o
critério que levou \citet{neto2022kw} a escolher o SQLite quando o banco foi
criado.

O esquema trata cada uso do kw como um evento ocorrido numa data e
hora \citep{neto2022kw}. A tabela \texttt{event} guarda essa data e hora, e a
tabela \texttt{executed} a estende, com a mesma chave, para dizer qual
operação rodou, quanto tempo levou e em que estado terminou; os nomes das
operações e dos estados ficam em tabelas próprias
(Figura~\ref{fig:esquema-kw}). A mesma tabela \texttt{event} serve de base
para as sessões do cronômetro de foco, que não entram aqui. O código do kw
não escreve nessas tabelas diretamente. A função \texttt{statistics\_manager}
insere uma linha plana na \textit{view} \texttt{statistics\_report}, e um
gatilho \texttt{INSTEAD OF INSERT} distribui o dado entre as tabelas
\citep{kworkflowrepo}. A mesma \textit{view} é por onde o \texttt{kw report}
lê, e as suas seis colunas (identificador, operação, estado, data, hora e
duração) são tudo o que se sabe de cada execução.

\begin{figure}[htbp]
  \centering
  \resizebox{\textwidth}{!}{%!TeX root=../main.tex
% Figura 2.4: parte do esquema do banco do kworkflow usada neste trabalho.
% Desenhada a partir de database/kwdb.sql (commit 2d8a04b). Mostra as tabelas
% do registro de execuções e a view statistics_report, que é a porta de
% entrada (gatilho INSTEAD OF INSERT) e de saída (kw report) desses dados.
% As tabelas do cronômetro de foco (timebox, tag) aparecem só pelo nome.
\begin{tikzpicture}[
    x=1cm, y=1cm,
    every node/.style={font=\footnotesize},
    tabela/.style={draw, rectangle split, rectangle split parts=2,
                   rectangle split part fill={black!12, white},
                   inner sep=3pt, font=\footnotesize},
    fora/.style={draw, dashed, text=black!55, draw=black!40,
                 inner sep=3pt, font=\footnotesize\ttfamily},
    visao/.style={draw, very thick, rounded corners=4pt, fill=black!4,
                  inner sep=4pt, align=center, font=\footnotesize},
    fk/.style={-{Latex[length=1.6mm]}, semithick},
    liga/.style={semithick, dashed},
  ]
  % Colunas: nome em \texttt à esquerda, tipo ou chave à direita
  \providecommand{\kwcolunas}{}\renewcommand{\kwcolunas}[1]{%
    \begin{tabular}{@{}l@{\hspace{1.2em}}r@{}}#1\end{tabular}}
  \providecommand{\kwcol}{}\renewcommand{\kwcol}[2]{\texttt{#1} & {\scriptsize #2}\\}

  \node[tabela] (event) at (0, 0) {%
    \texttt{event}
    \nodepart{two}
    \kwcolunas{\kwcol{id}{PK}\kwcol{date}{TEXT}\kwcol{time}{TEXT}}};

  \node[tabela] (executed) at (5.1, 0) {%
    \texttt{executed}
    \nodepart{two}
    \kwcolunas{\kwcol{id}{PK, FK}\kwcol{command\_label\_id}{FK}\kwcol{status\_id}{FK}%
             \kwcol{elapsed\_time\_in\_secs}{INTEGER}}};

  \node[tabela] (label) at (9.8, 0.8) {%
    \texttt{command\_label}
    \nodepart{two}
    \kwcolunas{\kwcol{id}{PK}\kwcol{name}{TEXT}}};

  \node[tabela] (status) at (9.8, -0.95) {%
    \texttt{status}
    \nodepart{two}
    \kwcolunas{\kwcol{id}{PK}\kwcol{name}{TEXT}}};

  \node[fora] (timebox) at (0, -2.1) {timebox, tag};

  % Chaves estrangeiras
  \draw[fk] (executed.west) -- node[above, font=\scriptsize] {mesmo \texttt{id}} (event.east);
  \draw[fk] ([yshift=1.5mm]executed.east) -- (label.west);
  \draw[fk] ([yshift=-2.5mm]executed.east) -- (status.west);
  \draw[fk, draw=black!40] (timebox.north) -- (event.south);

  % A view por cima
  \node[visao] (view) at (5.1, 2.6) {%
    view \texttt{statistics\_report}\\[1pt]
    {\ttfamily id, label\_name, status, date, time, elapsed\_time\_in\_secs}};
  \draw[liga] ([xshift=4mm]view.south west) -- (event.north);
  \draw[liga] (view.south) -- (executed.north);
  \draw[liga] ([xshift=-4mm]view.south east) -- (label.north);
\end{tikzpicture}
}
  \caption{Parte do esquema do \texttt{kw.db} que guarda o registro de
    execuções, com a \textit{view} usada para escrever e ler os dados.
    Desenhado a partir de \texttt{database/kwdb.sql} \citep{kworkflowrepo}.}
  \label{fig:esquema-kw}
\end{figure}

O design serve bem ao que foi feito para servir: o esquema é normalizado,
fácil de consultar e tem um único ponto de escrita. Três propriedades dele,
no entanto, só fazem diferença quando os dados precisam sair da máquina, e as
três voltam no Capítulo~\ref{cap:solucao}.

A primeira é que o identificador
de cada evento cresce na ordem de inserção, o que permite saber o que já foi
lido; como a coluna não é declarada com \texttt{AUTOINCREMENT}, o SQLite pode
reaproveitar o identificador de um evento apagado, embora nenhum código do kw
apague eventos hoje. A segunda é que data e hora são guardadas como texto, em
colunas separadas, no horário local e sem fuso. A terceira é que o esquema não
tem número de versão: todas as tabelas são criadas com \texttt{CREATE TABLE IF
NOT EXISTS}, e o campo \texttt{user\_version} do SQLite fica em zero
\citep{kworkflowrepo}. Quando o esquema foi desenhado, em 2021, os dados
nunca deixavam a máquina de quem os gerou, e as três escolhas eram as
óbvias. Elas continuam corretas para o \texttt{kw report}. O que não existe
ainda é um caminho para que esses dados saiam dali.
